% ============================================================================
%  Scale-Packing Residual Analysis — LAMMPS Validation Layer
%  VSEPR-SIM-ORGANICS & METALLICS {release}
%  99_backmatter — Complete Back Matter
%  Appendices A–F, Glossary, References, Future Work
%  Compile with pdflatex (twice for refs).
% ============================================================================
\documentclass[11pt, a4paper, oneside]{article}

% --- Geometry ----------------------------------------------------------------
\usepackage[
  top=1.0in, bottom=1.0in, left=1.15in, right=1.15in,
  headheight=14pt
]{geometry}

% --- Fonts / encoding --------------------------------------------------------
\usepackage[T1]{fontenc}
\usepackage{lmodern}
\usepackage{microtype}

% --- Math --------------------------------------------------------------------
\usepackage{amsmath, amssymb, amsthm, mathtools}
\usepackage{bm}

% --- Tables ------------------------------------------------------------------
\usepackage{booktabs}
\usepackage{array}
\usepackage{tabularx}
\usepackage{longtable}
\usepackage{multirow}

% --- Code listings -----------------------------------------------------------
\usepackage{listings}
\usepackage{xcolor}
\definecolor{codegray}{rgb}{0.5,0.5,0.5}
\definecolor{codeblue}{rgb}{0.13,0.13,0.9}
\definecolor{codegreen}{rgb}{0.0,0.5,0.0}
\definecolor{codebg}{rgb}{0.97,0.97,0.97}
\lstset{
  backgroundcolor=\color{codebg},
  basicstyle=\ttfamily\footnotesize,
  keywordstyle=\color{codeblue}\bfseries,
  commentstyle=\color{codegreen}\itshape,
  numberstyle=\tiny\color{codegray},
  numbers=left, numbersep=6pt,
  breaklines=true, frame=single,
  rulecolor=\color{codegray},
  captionpos=b, showstringspaces=false,
  tabsize=2
}

% --- Lists -------------------------------------------------------------------
\usepackage{enumitem}
\setlist{nosep, leftmargin=1.5em}

% --- Headers / footers -------------------------------------------------------
\usepackage{fancyhdr}
\pagestyle{fancy}
\fancyhf{}
\fancyhead[L]{\small\textit{Scale-Packing MD Implementation}}
\fancyhead[R]{\small\textit{Back Matter}}
\fancyfoot[C]{\thepage}
\renewcommand{\headrulewidth}{0.4pt}

% --- Section formatting ------------------------------------------------------
\usepackage{titlesec}
\titleformat{\section}{\Large\bfseries}{\S\thesection}{1em}{}
\titleformat{\subsection}{\large\bfseries}{\thesubsection}{1em}{}
\titleformat{\subsubsection}{\normalsize\bfseries}{\thesubsubsection}{1em}{}

% --- Hyperlinks --------------------------------------------------------------
\usepackage[colorlinks=true, linkcolor=black, citecolor=black, urlcolor=blue]{hyperref}

% --- Appendix numbering ------------------------------------------------------
\usepackage{appendix}

% --- Custom operators --------------------------------------------------------
\newcommand{\Pack}[1]{\Pi_{#1}}
\newcommand{\Ann}{\mathcal{A}}
\newcommand{\Comp}[1]{\mathbf{C}_{#1}}
\newcommand{\Ident}[1]{\mathbf{I}_{#1}}
\newcommand{\Resid}[1]{\mathbf{R}_{#1}}
\newcommand{\Proj}[2]{P_{#1 \leftarrow #2}}
\newcommand{\Embed}[2]{E_{#2 \leftarrow #1}}
\newcommand{\ProjEM}{P_{\mathrm{EM}}}
\newcommand{\ProjG}{P_{G}}
\newcommand{\ProjDet}{P_{\mathrm{det}}}
\newcommand{\RecDet}{\mathcal{R}_{\mathrm{det}}}
\newcommand{\TrueLoc}{\xi_i^*}
\newcommand{\ObsLoc}{r_i^{\mathrm{obs}}}
\newcommand{\FuzzR}{r_f}
\newcommand{\ScaleL}{\ell_S}
\newcommand{\FuzzRatio}{\Phi_{i,S}}
\newcommand{\HiddenAct}{H_j}
\newcommand{\DiagMode}{\Lambda}
\newcommand{\DarkDens}{\rho_{\mathrm{dark}}}

% ============================================================================
\begin{document}

\begin{titlepage}
  \centering
  \vspace*{1.5cm}
  {\Huge\bfseries Scale-Packing Residual Analysis\par}
  \vspace{0.5cm}
  {\LARGE LAMMPS Validation Layer\par}
  \vspace{0.3cm}
  {\large VSEPR-SIM-ORGANICS \& METALLICS \{release\}\par}
  \vspace{1.5cm}
  {\LARGE\bfseries Back Matter\par}
  \vspace{0.5cm}
  {\large Appendices A--F\,|\,Glossary\,|\,References\,|\,Future Work\par}
  \vspace{2cm}
  {\normalsize Verification is the focus.\par}
  \vspace{1cm}
  {\large VSEPR-SIM Theoretical Division\par}
  {\large Development Day 64\par}
\end{titlepage}

\tableofcontents
\newpage

% ============================================================================
\begin{appendices}

% ============================================================================
\section{Notation}
\label{app:A}
% ============================================================================

This appendix centralises every symbol used in the implementation document.
All entries are listed in definition order, not alphabetical order, so that
the logical construction of the framework is visible in the table.

\subsection*{A.1 Core Packing Algebra}

\begin{longtable}{p{3.2cm} p{2.2cm} p{8.8cm}}
  \toprule
  \textbf{Symbol} & \textbf{Type} & \textbf{Definition} \\
  \midrule
  \endhead
  $\Pack{n}$ & operator &
	Packing operator at scale $n$. Maps component matrix
	$\Comp{n}$ to packed identity vector $\Ident{n+1}$:
	$\Comp{n} \xrightarrow{\Pack{n}} \Ident{n+1}$.
	Packing is irreversible without energy $\geq E_{\mathrm{bind}}$. \\[4pt]

  $\Ann$ & operator &
	Annihilation / unpacking map.
	$\Ident{n+1} \xrightarrow{\Ann} \Comp{n,\mathrm{out}}$.
	Destroys or transforms the original packed identity.
	Used in CERN-layer collision analysis. \\[4pt]

  $\Comp{n}$ & matrix &
	Component matrix at scale $n$.
	Rows are components; columns are physical channels
	(charge, mass, position, velocity, \ldots).
	Populated from atomistic trajectory fields. \\[4pt]

  $\Ident{n+1}$ & vector &
	Packed identity vector at scale $n+1$.
	Result of applying $\Pack{n}$ to $\Comp{n}$.
	For chemistry: molecule-level descriptor.
	For solids: unit-cell or grain-level descriptor. \\[4pt]

  $\Resid{n}$ & vector / tensor &
	Packing residual at scale $n$:
	$\Resid{n} = \Comp{n} - \Embed{n}{n+1}\!\bigl(\Ident{n+1}\bigr)$.
	Non-zero when packing is lossy. \\[4pt]

  \midrule
  \multicolumn{3}{l}{\textit{Projection and embedding}} \\
  \midrule

  $\Proj{A}{B}$ & operator &
	Projection from space $B$ into space $A$.
	$\Proj{C}{A}$: atomistic $\to$ chemical.
	$\Proj{S}{A}$: atomistic $\to$ solid-state.
	$\Proj{R}{A}$: atomistic $\to$ residual descriptor. \\[4pt]

  $\Embed{A}{B}$ & operator &
	Embedding / reconstruction from $A$ back into $B$.
	$\Embed{C}{A}$: chemical identity $\to$ atomistic-space prediction.
	Used to measure projection loss. \\[4pt]

  $\ProjEM$ & operator &
	Electromagnetic / visible projection.
	Retains only components at scales with
	$L_m \leq L_{\mathrm{EM}}$. \\[4pt]

  $\ProjG$ & operator &
	Gravitational projection.
	Couples to all scale levels; $\int K_G(w)\,dw = 1$. \\[4pt]

  $\ProjDet$ & operator &
	Detector forward projection.
	Maps parent collision identity to expected detector signals.
	Reserved for CERN high-energy verification layer. \\[4pt]

  $\RecDet$ & operator &
	Detector reconstruction operator.
	Inverse of $\ProjDet$ over observed final-state debris.
	Reserved for CERN layer. \\[4pt]

  \midrule
  \multicolumn{3}{l}{\textit{Position and fuzz metrics}} \\
  \midrule

  $\TrueLoc$ & vector &
	True identity-location barycenter (identity anchor) of particle $i$:
	$\TrueLoc = \frac{\sum_k \alpha_k c_k}{\sum_k \alpha_k}$.
	Not directly observable. \\[4pt]

  $\ObsLoc$ & vector &
	Observed 3D projected position of particle $i$:
	$\ObsLoc = P_{3D}(\TrueLoc)$.
	The only quantity accessible to the MD dump. \\[4pt]

  $\FuzzR$ & scalar &
	Fuzz radius: weighted RMS spread of component 3D projections
	around the observed position:
	$\FuzzR^2 = \frac{\sum_k \alpha_k \|P_{3D}(c_k) - \ObsLoc\|^2}
					 {\sum_k \alpha_k}$. \\[4pt]

  $\ScaleL$ & scalar &
	Scale resolution: characteristic length of scale $S$.
	Used to normalise the fuzz radius. \\[4pt]

  $\FuzzRatio$ & scalar &
	Fuzz ratio: $\FuzzRatio = \FuzzR / \ScaleL$.
	$\FuzzRatio \ll 1$: point-like at scale $S$.
	$\FuzzRatio \sim 1$: fuzz-zone.
	$\FuzzRatio > 1$: non-point-defined at scale $S$. \\[4pt]

  \midrule
  \multicolumn{3}{l}{\textit{Hidden activity and residual modes}} \\
  \midrule

  $\HiddenAct$ & scalar &
	Hidden column activity for component $j$:
	$H_j = \mathbf{c}_j^\top \mathbf{W}\, \mathbf{c}_j$.
	Non-zero even when the observable channel
	$I_j = \bm{\alpha}^\top \mathbf{c}_j = 0$. \\[4pt]

  $\DiagMode$ & matrix &
	Diagonalised internal mode intensity matrix.
	Eigenvalues of $\mathbf{C}_n^\top \mathbf{W} \mathbf{C}_n$.
	Each non-zero eigenvalue represents a structurally active
	mode invisible in the projected observable. \\[4pt]

  $\DarkDens$ & scalar field &
	Gravitational-minus-visible projection residual density:
	$\DarkDens(x,y,z;\,t) = \int \rho(x,y,z,w;\,t)
	[K_G(w) - K_{\mathrm{EM}}(w)]\,dw$.
	Defined for completeness; not computed in LAMMPS layer. \\[4pt]

  \bottomrule
\end{longtable}

\subsection*{A.2 Index conventions}

\begin{tabular}{ll}
  \toprule
  Index & Meaning \\
  \midrule
  $i$ & Particle / atom index (1-indexed, matching LAMMPS \texttt{id}) \\
  $j$ & Component channel index within a component matrix column \\
  $k$ & Component row index within a component matrix \\
  $n$ & Scale level (0 = existence, 1--5 standard, $D$ = dark) \\
  $t$ & Trajectory frame time (LAMMPS timestep units) \\
  $f$ & Frame index in trajectory tensor \\
  \bottomrule
\end{tabular}

\subsection*{A.3 Typeface conventions}

\begin{tabular}{ll}
  \toprule
  Typeface & Meaning \\
  \midrule
  $\mathbf{C}$ & Component matrix (bold, capital) \\
  $\mathbf{I}$ & Packed identity (bold, capital; not identity matrix) \\
  $\mathbf{R}$ & Residual vector or tensor (bold, capital) \\
  $\mathbf{W}$ & Hidden-intensity weight matrix (bold, capital) \\
  $\bm{\alpha}$ & Projection weight vector (bold, lowercase Greek) \\
  $P, E$ & Projection and embedding operators (italic capital) \\
  $\mathcal{A}$ & Annihilation map (calligraphic) \\
  $\mathcal{I}$ & Full identity state (calligraphic, when not matrix) \\
  \bottomrule
\end{tabular}

% ============================================================================
\section{LAMMPS Field Map}
\label{app:B}
% ============================================================================

This appendix defines exactly which LAMMPS \texttt{dump custom} fields are
required, preferred, or optional for the atomistic component matrix, and
specifies the handling rule for each.

\subsection*{B.1 Minimum required dump command}

\begin{lstlisting}[language=bash, caption={Minimum LAMMPS dump for
	Scale-Packing analysis}]
dump    spra all custom 100 trajectory.dump &
		id type mol x y z vx vy vz fx fy fz q mass
dump_modify spra sort id
\end{lstlisting}

If \texttt{mol} is not defined by the atom style, replace with \texttt{0}
via a constant column or omit and flag as missing.
If \texttt{q} is not defined, set column to \texttt{0.0} and record
\texttt{charge\_available = false} in the metadata block.

\subsection*{B.2 Complete field map}

\begin{longtable}{p{2.2cm} p{2.5cm} p{1.8cm} p{7.1cm}}
  \toprule
  \textbf{LAMMPS field} & \textbf{Component matrix column} &
  \textbf{Status} & \textbf{Notes} \\
  \midrule
  \endhead

  \texttt{id}       & \texttt{atom\_id}     & Required &
	1-indexed integer. Row identity for the component matrix. \\

  \texttt{type}     & \texttt{atom\_type}   & Required &
	Integer type index. Mapped to element symbol via
	LAMMPS \texttt{mass} table or user-supplied type map. \\

  \texttt{mol}      & \texttt{mol\_id}      & Preferred &
	Molecule-group ID. Used to assemble $\Comp{A}$ rows into
	molecule-level packing groups. Set to 0 if atom style has
	no molecule concept (solids). \\

  \texttt{x y z}    & \texttt{x, y, z}     & Required &
	Wrapped coordinates. Used when \texttt{xu yu zu} not available.
	Periodic-unwrapping applied using \texttt{ix iy iz} image flags. \\

  \texttt{xu yu zu} & \texttt{x, y, z}     & Preferred &
	Unwrapped coordinates. Use in preference to wrapped
	\texttt{x y z} when available. Bypasses image-flag arithmetic. \\

  \texttt{vx vy vz} & \texttt{vx, vy, vz}  & Required &
	Per-atom velocity. Used for kinetic energy and momentum
	channels in $\Comp{A}$. \\

  \texttt{fx fy fz} & \texttt{fx, fy, fz}  & Required &
	Per-atom force. Used for force-residual and stress-proxy
	channels. \\

  \texttt{q}        & \texttt{charge}       & Preferred &
	Per-atom charge. Required for charge-projection operators
	and hidden column activity in ionic systems.
	Set to 0.0 if not available; flag \texttt{charge\_available}
	accordingly. \\

  \texttt{mass}     & \texttt{mass}         & Required &
	Per-atom mass (a.m.u.). Used in anchor weight $\alpha_k$
	and fuzz-radius computation. \\

  \texttt{ix iy iz} & (internal)           & Preferred &
	Image flags for periodic-boundary unwrapping.
	Required when using wrapped \texttt{x y z}. \\

  \texttt{timestep} & \texttt{frame\_time}  & Required &
	Simulation time at dump frame. Stored in trajectory metadata;
	not a component matrix column. \\

  Box bounds        & \texttt{cell}         & Required &
	Simulation cell vectors. Read from dump header.
	Required for periodic minimum-image distance. \\

  \midrule

  \texttt{c\_pe[]}  & \texttt{pe\_atom}     & Optional &
	Per-atom potential energy (requires \texttt{compute pe/atom}).
	Used as additional energy channel when available. \\

  \texttt{c\_stress[][]} & \texttt{stress}  & Optional &
	Per-atom stress tensor components.
	Used in solid-state strain-residual operators. \\

  \texttt{c\_cn[]}  & \texttt{coord\_num}   & Optional &
	Per-atom coordination number (requires neighbour compute).
	Used in solid-state neighbor graph construction. \\

  \bottomrule
\end{longtable}

\subsection*{B.3 Missing-field policy}

\begin{enumerate}
  \item \textbf{Required field absent}: abort parsing; emit error
	\texttt{[FIELD\_MISSING] <field>} to analysis log.
  \item \textbf{Preferred field absent}: continue; set column to
	\texttt{0.0}; write \texttt{[FIELD\_DEFAULT] <field> = 0.0} to log.
  \item \textbf{Optional field absent}: silently skip dependent
	operators; mark operator as \texttt{INACTIVE} in metadata.
  \item \textbf{Wrapped vs.\ unwrapped coordinates}: if
	\texttt{xu yu zu} absent, apply unwrapping:
	\[
	  x_i^{\mathrm{unwrap}} = x_i + i_x \cdot L_x, \quad
	  y_i^{\mathrm{unwrap}} = y_i + i_y \cdot L_y, \quad
	  z_i^{\mathrm{unwrap}} = z_i + i_z \cdot L_z
	\]
	where $L_x, L_y, L_z$ are box lengths and
	$i_x, i_y, i_z$ are image flags.
	If image flags also absent: flag \texttt{unwrap\_unavailable}.
\end{enumerate}

% ============================================================================
\section{Projection Operator Catalog}
\label{app:C}
% ============================================================================

Each operator entry states: input space, output space, recovered variables,
lost variables, residual definition, metadata required, and validation method.

\subsection*{C.1 Operator naming convention}

\[
  P_{X \leftarrow Y}: \text{ projection from space } Y \text{ into space } X
\]

Operators are composable:
$P_{R \leftarrow A} \circ P_{A \leftarrow M}$ maps molecule-level to
residual space via atomistic intermediate.

\subsection*{C.2 Catalog}

\subsubsection*{$P_{C \leftarrow A}$ --- Atomistic to Chemical}

\begin{tabularx}{\linewidth}{l X}
  \toprule
  Input space & Atomistic component matrix $\Comp{A}$
	(one row per atom; columns: position, velocity, force,
	charge, mass) \\
  Output space & Chemical identity vector $\Ident{C}$
	(one row per molecule; columns: molecular formula,
	center of mass, total charge, total mass, bond graph hash) \\
  Recovered & Molecular identity, stoichiometry, charge, CoM, geometry \\
  Lost & Per-atom sub-molecular detail; intra-molecular
	velocity spread; internal force distribution \\
  Residual & $\Resid{A \to C} = \Comp{A}
	- \Embed{C}{A}\!\bigl(\Ident{C}\bigr)$. Per-atom deviation
	from the chemical-identity prediction. \\
  Metadata & Molecule group assignments (\texttt{mol\_id});
	bond graph at current frame; atom-type map \\
  Validation & NaCl crystal: $Q_\text{mol} = 0$ for each NaCl pair;
	water cluster: $Q_\text{mol} = 0$, CoM stable \\
  \bottomrule
\end{tabularx}

\vspace{1em}

\subsubsection*{$P_{S \leftarrow A}$ --- Atomistic to Solid-State}

\begin{tabularx}{\linewidth}{l X}
  \toprule
  Input space & Atomistic component matrix $\Comp{A}$ \\
  Output space & Solid-state identity vector $\Ident{S}$
	(one row per unit cell or grain;
	columns: lattice parameter estimates, packing fraction,
	mean coordination number, defect density) \\
  Recovered & Crystal symmetry proxies, packing order, neighbour structure \\
  Lost & Individual atomic displacements from ideal lattice sites;
	local stress heterogeneity \\
  Residual & $\Resid{A \to S}$: per-atom deviation from
	ideal lattice prediction. Large values indicate
	defects, strain, or amorphization. \\
  Metadata & Reference lattice parameters; neighbour cutoff radius;
	species type map \\
  Validation & Silicon lattice: residual near zero at equilibrium;
	spike at vacancy site \\
  \bottomrule
\end{tabularx}

\vspace{1em}

\subsubsection*{$P_{R \leftarrow A}$ --- Atomistic to Residual Descriptor}

\begin{tabularx}{\linewidth}{l X}
  \toprule
  Input space & Atomistic component matrix $\Comp{A}$ \\
  Output space & Residual descriptor vector per atom:
	(fuzz radius, hidden column activity,
	diagonal mode intensity, packing loss norm) \\
  Recovered & Structured diagnostic: where projection fails and how \\
  Lost & Not applicable; this is a loss-measuring operator \\
  Residual & This operator \emph{is} the residual output. No further
	residual is defined. \\
  Metadata & Scale resolution $\ScaleL$ for fuzz-ratio normalisation;
	weight matrix $\mathbf{W}$ per channel \\
  Validation & Argon LJ gas: fuzz radius $\sim k_BT$ thermal spread;
	NaCl: hidden charge activity $> 0$ despite neutral pairs \\
  \bottomrule
\end{tabularx}

\vspace{1em}

\subsubsection*{$P_{R_\mathrm{chem} \leftarrow A}$ --- Reactive Chemical Residual}

\begin{tabularx}{\linewidth}{l X}
  \toprule
  Input space & Atomistic component matrix + bond-order matrix $\mathbf{B}$
	at frames $t$ and $t + \Delta t$ \\
  Output space & Reactive event residual: bond-change delta,
	charge redistribution, species-transition confidence \\
  Recovered & Bond formation/breaking events; charge flow; species
	identity change \\
  Lost & Sub-timestep bond dynamics; quantum transition amplitudes \\
  Residual & $\Resid{\mathrm{react}}$: difference between
	graph state before and after event; includes bond-order
	fractional change and partial-charge delta \\
  Metadata & Bond-order cutoffs; species definition table;
	persistence threshold (number of frames) \\
  Validation & ReaxFF hydrocarbon case: bond-breaking events match
	known reaction channels \\
  \bottomrule
\end{tabularx}

\vspace{1em}

\subsubsection*{$P_{\mathrm{EM}}$ --- Electromagnetic Projection}

\begin{tabularx}{\linewidth}{l X}
  \toprule
  Input space & Full identity state $\mathcal{I}$ over scale ladder \\
  Output space & EM-observable density:
	retains components with $L_m \leq L_{\mathrm{EM}}$ \\
  Recovered & Charge distribution, photon-coupled interactions \\
  Lost & Higher-scale components ($w > L_{\mathrm{EM}}$) \\
  Residual & $\DarkDens = P_G(\rho) - P_{\mathrm{EM}}(\rho)$ \\
  Notes & Used in theoretical layer only; not computed in LAMMPS kernel \\
  \bottomrule
\end{tabularx}

\vspace{1em}

\subsubsection*{$P_{\mathrm{det}}$ and $\mathcal{R}_{\mathrm{det}}$ --- Detector Projection / Reconstruction}

\begin{tabularx}{\linewidth}{l X}
  \toprule
  Status & \textbf{Reserved.} CERN high-energy verification layer only.
	Not implemented in LAMMPS-primary release. \\
  Purpose & $P_{\mathrm{det}}$: maps collision parent identity to
	expected detector hit pattern.
	$\mathcal{R}_{\mathrm{det}}$: reconstructs parent identity
	from observed final-state debris. \\
  Validation target & $J/\psi \to \mu^+\mu^-$;
	$H \to ZZ \to 4\ell$ channels \\
  \bottomrule
\end{tabularx}

% ============================================================================
\section{Residual Metric Definitions}
\label{app:D}
% ============================================================================

This appendix defines every residual metric to ensure that outputs across
chemistry, solids, and reactive cases are numerically comparable.

\subsection*{D.1 Packing loss norms}

\textbf{Chemical packing loss:}
\begin{equation}
  \varepsilon_{A \to C}
  = \frac{\|\Resid{A \to C}\|_F}{\|\Comp{A}\|_F}
  \tag{D.1}
\end{equation}

\textbf{Solid packing loss:}
\begin{equation}
  \varepsilon_{A \to S}
  = \frac{\|\Resid{A \to S}\|_F}{\|\Comp{A}\|_F}
  \tag{D.2}
\end{equation}

\textbf{Reactive packing loss:}
\begin{equation}
  \varepsilon_{\mathrm{react}}
  = \frac{\|\Resid{\mathrm{react}}\|_1}{N_{\mathrm{atoms,event}}}
  \tag{D.3}
\end{equation}

where $\|\cdot\|_F$ is the Frobenius norm and $\|\cdot\|_1$ is the
element-wise $\ell^1$ norm. Normalisation by $\|\Comp{A}\|_F$ makes
$\varepsilon \in [0, 1]$.

\subsection*{D.2 Residual vectors}

\textbf{Atomistic-to-chemical residual} (per atom $i$):
\begin{equation}
  \mathbf{R}_{i,\,A \to C}
  = \mathbf{c}_i^{(A)} - \Embed{C}{A}\!\bigl(\Ident{C}[\mathrm{mol}(i)]\bigr)
  \tag{D.4}
\end{equation}

\textbf{Atomistic-to-solid residual} (per atom $i$):
\begin{equation}
  \mathbf{R}_{i,\,A \to S}
  = \mathbf{c}_i^{(A)} - \mathbf{c}_i^{\mathrm{lattice}}
  \tag{D.5}
\end{equation}
where $\mathbf{c}_i^{\mathrm{lattice}}$ is the atom's predicted position
and velocity on the ideal reference lattice.

\subsection*{D.3 Hidden column activity}

For component vector $\mathbf{c}_j \in \mathbb{R}^m$:
\begin{equation}
  I_j = \bm{\alpha}^\top \mathbf{c}_j
  \tag{D.6}
\end{equation}
\begin{equation}
  H_j = \mathbf{c}_j^\top \mathbf{W}\, \mathbf{c}_j
  \tag{D.7}
\end{equation}

Hidden-active condition: $I_j = 0$, $H_j > 0$.
Standard choice $\mathbf{W} = \mathbf{I}$ gives $H_j = \|\mathbf{c}_j\|^2$.

\subsection*{D.4 Diagonal residual mode intensity}

Let $\mathbf{G} = \Comp{n}^\top \mathbf{W}\, \Comp{n}$. Eigendecompose:
\begin{equation}
  \mathbf{G} = \mathbf{Q}\, \DiagMode\, \mathbf{Q}^\top
  \tag{D.8}
\end{equation}
$\DiagMode = \mathrm{diag}(\lambda_1, \lambda_2, \ldots, \lambda_m)$.
Each $\lambda_j > 0$ indicates an internal mode with non-zero intensity
even when the corresponding observable channel sums to zero.

\subsection*{D.5 Fuzz radius and fuzz ratio}

\begin{equation}
  \FuzzR^2
  = \frac{\sum_k \alpha_k
		  \bigl\|P_{3D}(\mathbf{c}_k) - \ObsLoc\bigr\|^2}
		 {\sum_k \alpha_k}
  \tag{D.9}
\end{equation}
\begin{equation}
  \FuzzRatio = \frac{\FuzzR}{\ScaleL}
  \tag{D.10}
\end{equation}

\subsection*{D.6 Threshold bands}

\begin{center}
\begin{tabular}{lll}
  \toprule
  Metric & Healthy range & Elevated / investigate \\
  \midrule
  $\varepsilon_{A \to C}$ & $< 0.05$ & $> 0.15$ (reactive transition) \\
  $\varepsilon_{A \to S}$ & $< 0.02$ at $T=0$ & $> 0.10$ (defect / melt) \\
  $\varepsilon_{\mathrm{react}}$ & $< 0.01$ (no event) & $> 0.05$ (event) \\
  $H_j / \|\mathbf{c}_j\|^2$ & $< 0.1$ (charge-neutral) & $> 0.5$ (internal activity) \\
  $\FuzzRatio$ & $\ll 1$ (point-like) & $\sim 1$ (fuzz-zone) \\
  \bottomrule
\end{tabular}
\end{center}

\textit{Threshold values are initial estimates subject to calibration
against LAMMPS validation cases in Section~10 of the main document.}

% ============================================================================
\section{Reactive Event Schema}
\label{app:E}
% ============================================================================

\subsection*{E.1 JSON event record format}

Each reactive event is recorded as one JSON object.
The full trajectory event log is a JSON array of these objects.

\begin{lstlisting}[language={},caption={Reactive event record (JSON schema)}]
{
  "event_id":          "string (UUID-4)",
  "time":              "float  (LAMMPS timestep units)",
  "frame":             "int    (trajectory frame index, 0-based)",
  "event_type":        "string (see E.2)",
  "atoms_involved":    ["int", "..."],
  "bonds_formed":      [["int","int"], "..."],
  "bonds_broken":      [["int","int"], "..."],
  "charge_delta":      "float  (net charge change summed over atoms_involved)",
  "species_before":    {"atom_id": "symbol", "..."},
  "species_after":     {"atom_id": "symbol", "..."},
  "graph_before_hash": "string (FNV-1a of adjacency list before event)",
  "graph_after_hash":  "string (FNV-1a of adjacency list after event)",
  "residual_before":   "float  (||R_react|| before event)",
  "residual_after":    "float  (||R_react|| after event)",
  "confidence_score":  "float  in [0, 1]  (see E.3)"
}
\end{lstlisting}

\subsection*{E.2 Event type enumeration}

\begin{longtable}{p{5.5cm} p{9.0cm}}
  \toprule
  \textbf{Event type string} & \textbf{Meaning} \\
  \midrule
  \endhead
  \texttt{bond\_formation}        & New bond detected between two atoms whose
	bond order crossed threshold from below. \\
  \texttt{bond\_breaking}         & Existing bond lost; bond order fell below
	threshold. \\
  \texttt{charge\_redistribution} & Per-atom charge changed by $> \delta q$
	without bond-order crossing. \\
  \texttt{species\_change}        & Atom identity (effective element or
	functional group) changed based on coordination
	and charge context. \\
  \texttt{coordination\_change}   & Coordination number changed by $\geq 1$
	without a bond-order threshold event. \\
  \texttt{fragmentation}          & Connected component of the molecular
	graph split into two or more disconnected graphs. \\
  \texttt{recombination}          & Two previously disconnected graph
	components merged into one. \\
  \texttt{ambiguous\_transition}  & Bond order is in the hysteresis window
	$[B_{\mathrm{lo}}, B_{\mathrm{hi}}]$ for $> N_{\mathrm{persist}}$
	frames. Not yet classified as formed or broken. \\
  \bottomrule
\end{longtable}

\subsection*{E.3 Confidence score definition}

\begin{equation}
  \mathrm{confidence} = w_b\,s_b + w_q\,s_q + w_p\,s_p
  \tag{E.1}
\end{equation}

where:
\begin{itemize}
  \item $s_b = \min(1,\, (B - B_{\mathrm{hi}}) / B_{\mathrm{range}})$:
	bond-order margin above hysteresis window
  \item $s_q = \min(1,\, |\Delta q| / \Delta q_{\mathrm{ref}})$:
	charge-delta relative to reference scale
  \item $s_p = N_{\mathrm{persist}} / N_{\mathrm{persist,\,ref}}$:
	persistence fraction
  \item Default weights: $w_b = 0.5$, $w_q = 0.3$, $w_p = 0.2$
\end{itemize}

Events with $\mathrm{confidence} < 0.3$ are tagged
\texttt{ambiguous\_transition} regardless of bond-order classification.

\subsection*{E.4 Persistence and hysteresis rule}

A bond-order change is not recorded as an event until it persists for
$N_{\mathrm{persist}} \geq 3$ consecutive frames. This suppresses
thermal-fluctuation noise at the bond-order threshold. The hysteresis
window $[B_{\mathrm{lo}}, B_{\mathrm{hi}}]$ should be set per
species pair based on the force field in use.

% ============================================================================
\section{Validation Assignment Table}
\label{app:F}
% ============================================================================

\subsection*{F.1 Validation case matrix}

\begin{longtable}{p{4.5cm} p{4.5cm} p{5.5cm}}
  \toprule
  \textbf{Validation case} & \textbf{Modules tested} & \textbf{Pass criterion} \\
  \midrule
  \endhead

  Argon LJ gas\newline
  (NVT, 300 K, 1000 atoms)
  &
  MD mapping, trajectory parsing, $P_{R \leftarrow A}$
  &
  Fuzz radius tracks $\sqrt{k_B T / m_{\mathrm{Ar}}}$;
  $\varepsilon_{A \to C} < 0.05$;
  no hidden column activity (monatomic) \\
  \midrule

  Water cluster\newline
  (H$_2$O)$_{32}$, NPT 300 K
  &
  $P_{C \leftarrow A}$, chemical projection, geometry, CoM
  &
  Each molecule: $Q_{\mathrm{mol}} = 0 \pm 10^{-6}e$;
  CoM velocity $\sim$ Maxwell-Boltzmann;
  O--H bond lengths $1.0 \pm 0.05$\,\AA \\
  \midrule

  NaCl crystal\newline
  ($4 \times 4 \times 4$ supercell, 300 K)
  &
  $P_{C \leftarrow A}$, $P_{S \leftarrow A}$, charge channels,
  hidden activity
  &
  Ion pairs: $Q_{\mathrm{pair}} = 0$;
  $H_j > 0$ (internal charge activity);
  $\varepsilon_{A \to S} < 0.02$ at equilibrium \\
  \midrule

  Silicon lattice\newline
  (diamond cubic, 512 atoms, 300 K)
  &
  $P_{S \leftarrow A}$, neighbour graph, lattice projection
  &
  All atoms $\mathrm{CN} = 4$;
  $\varepsilon_{A \to S} < 0.01$;
  fuzz ratio $\ll 1$ \\
  \midrule

  Vacancy + interstitial\newline
  (Si lattice $+$ Frenkel pair)
  &
  $P_{S \leftarrow A}$, $P_{R \leftarrow A}$, defect descriptors
  &
  Vacancy atom: $\FuzzRatio > 1$;
  $\varepsilon_{A \to S} > 0.10$ at defect sites;
  interstitial: $\mathrm{CN} \neq 4$ \\
  \midrule

  Diffusion case\newline
  (Na in NaCl at 900 K)
  &
  MSD computation, $D_{\mathrm{eff}}$, trajectory tensor
  &
  $D_{\mathrm{Na}}$ within 20\,\% of experimental value;
  MSD linear regime identified \\
  \midrule

  Reactive hydrocarbon\newline
  (ReaxFF CH$_4$ combustion, 3000 K)
  &
  $P_{R_{\mathrm{chem}} \leftarrow A}$, bond-order matrix,
  reactive event detection
  &
  C--H bond-breaking events detected with
  confidence $> 0.5$;
  no ghost events in equilibrium run \\
  \midrule

  Oxidation / corrosion\newline
  (Fe slab + O$_2$, ReaxFF)
  &
  Reactive residual, charge redistribution,
  species change
  &
  Fe$^{2+}$/Fe$^{3+}$ species detected;
  $\varepsilon_{\mathrm{react}} > 0.05$ at oxide growth sites \\
  \midrule

  Thermal solid case\newline
  (NaCl heated to 1000 K)
  &
  Fuzz radius, displacement residual, solid projection
  &
  $\FuzzRatio$ increases monotonically with $T$;
  $\varepsilon_{A \to S}$ jumps at melting transition \\
  \midrule

  \bottomrule
\end{longtable}

\subsection*{F.2 Case input file status}

\begin{tabular}{ll}
  \toprule
  Case & Input file status \\
  \midrule
  Argon LJ gas            & TODO \\
  Water cluster           & TODO \\
  NaCl crystal            & TODO \\
  Silicon lattice         & TODO \\
  Vacancy + interstitial  & TODO \\
  Diffusion case          & TODO \\
  Reactive hydrocarbon    & TODO \\
  Oxidation / corrosion   & TODO \\
  Thermal solid           & TODO \\
  \bottomrule
\end{tabular}

\textit{Input files (LAMMPS data + in-script) to be added to
\texttt{10\_validation\_examples/} as each case is constructed.}

\end{appendices}

% ============================================================================
\section*{Glossary}
\addcontentsline{toc}{section}{Glossary}
\label{sec:glossary}
% ============================================================================

\begin{description}[style=nextline, leftmargin=3.5cm, labelwidth=3.2cm]

  \item[Scale packing]
	The process by which lower-scale components
	$\Comp{n}$ self-assemble into a higher-scale packed identity
	$\Ident{n+1}$ via packing operator $\Pack{n}$.
	Information not preserved by this map becomes the packing residual.

  \item[Packed identity]
	A compressed representation of a system at a given scale level.
	For chemistry: a molecule-level descriptor.
	For solids: a unit-cell or grain-level descriptor.
	A packed identity exists at a scale \emph{above} its components.

  \item[Component matrix]
	The matrix $\Comp{n}$ whose rows are individual components at
	scale $n$ and whose columns are physical channels
	(position, velocity, charge, mass, force, \ldots).
	Populated directly from LAMMPS trajectory fields.

  \item[Projection operator]
	An operator $P_{X \leftarrow Y}$ that maps a state in space $Y$
	to a representation in space $X$, discarding information not
	observable at scale $X$.

  \item[Annihilation map]
	The inverse of packing: $\Ann: \Ident{n+1} \to \Comp{n,\mathrm{out}}$.
	Physically realised in collision experiments.
	Destroys the original packed identity.

  \item[Embedding operator]
	The operator $E_{Y \leftarrow X}$ that lifts a representation
	from space $X$ back into space $Y$ for comparison.
	Used to define the residual:
	$\Resid{} = \Comp{n} - \Embed{C}{A}(\Ident{C})$.

  \item[Residual]
	The difference between the full component state and the
	embedding of its packed identity:
	$\Resid{} = \Comp{n} - \Embed{n}{n+1}(\Ident{n+1})$.
	Non-zero when packing is lossy.
	Interpreted as a physically useful diagnostic,
	not as numerical error.

  \item[Hidden column activity]
	The quantity $H_j = \mathbf{c}_j^\top \mathbf{W} \mathbf{c}_j$
	for a component vector $\mathbf{c}_j$.
	Measures internal structural intensity that is invisible
	when the observable channel $I_j = \bm{\alpha}^\top \mathbf{c}_j = 0$.

  \item[Diagonal residual mode]
	An eigenvalue $\lambda_j > 0$ of the matrix
	$\mathbf{G} = \Comp{n}^\top \mathbf{W} \Comp{n}$.
	Each positive eigenvalue represents an internal mode
	with non-zero intensity that projects to zero in the
	observable channel.

  \item[Fuzz radius]
	The weighted RMS spread of component 3D projections
	around the observed position $\ObsLoc$.
	Measures how non-point-like a packed identity is at its
	observed scale.

  \item[Fuzz zone]
	The regime $\FuzzRatio \sim 1$ where a particle or system
	cannot be meaningfully treated as a point at the current
	observation scale. Not to be confused with classical positional
	uncertainty alone.

  \item[Atomistic projection]
	A projection from the full atom-level LAMMPS state into
	a higher-level (chemical or solid-state) descriptor space.
	Implemented as one of the $P_{X \leftarrow A}$ operators.

  \item[Chemical graph]
	The bond graph derived from per-atom positions and bond-order
	thresholds. Nodes are atoms; edges are bonds with order $> B_{\mathrm{lo}}$.
	The basis of reactive chemistry support.

  \item[Solid-state projection]
	The operator $P_{S \leftarrow A}$ that maps atomistic states
	to lattice-level descriptors: coordination numbers, packing
	fraction, defect density, strain proxy.

  \item[Reactive residual]
	The residual $\Resid{\mathrm{react}}$ measuring the
	structural change in the molecular graph and charge state
	across a reactive event window.

  \item[Detector projection]
	The operator $P_{\mathrm{det}}$ mapping a collision parent
	identity to expected detector signals.
	Reserved for CERN high-energy verification layer.

  \item[Detector reconstruction]
	The inverse operator $\mathcal{R}_{\mathrm{det}}$ that
	reconstructs a parent identity estimate from observed
	final-state detector hits.
	Reserved for CERN layer.

  \item[LAMMPS validation layer]
	The primary verification stage of this implementation.
	Uses LAMMPS trajectories for chemistry- and solids-focused
	MD validation. Runs before CERN data is introduced.

  \item[CERN high-energy verification layer]
	A later verification stage involving advanced projection,
	reconstruction, annihilation mapping, and detector-scale
	identity separation. Not implemented in this release.

\end{description}

% ============================================================================
\section*{References}
\addcontentsline{toc}{section}{References}
\label{sec:references}
% ============================================================================

\subsection*{LAMMPS}

\begin{enumerate}[label={[\arabic*]}, leftmargin=2.2em]

  \item A.P.\ Thompson et al.,
	``LAMMPS --- A flexible simulation tool for particle-based
	materials modeling at the atomic, meso, and continuum scales,''
	\textit{Computer Physics Communications}, 271 (2022) 108171.
	\url{https://doi.org/10.1016/j.cpc.2021.108171}

  \item LAMMPS Documentation: \texttt{dump custom} command.
	\url{https://docs.lammps.org/dump.html}

  \item LAMMPS Documentation: ReaxFF pair style (\texttt{pair\_style reaxff}).
	\url{https://docs.lammps.org/pair_reaxff.html}

\end{enumerate}

\subsection*{Reactive force fields}

\begin{enumerate}[resume, label={[\arabic*]}, leftmargin=2.2em]

  \item A.C.T.\ van Duin et al.,
	``ReaxFF: A Reactive Force Field for Hydrocarbons,''
	\textit{Journal of Physical Chemistry A}, 105 (2001) 9396.

  \item K.\ Chenoweth, A.C.T.\ van Duin, W.A.\ Goddard III,
	``ReaxFF Reactive Force Field for Molecular Dynamics
	Simulations of Hydrocarbon Oxidation,''
	\textit{Journal of Physical Chemistry A}, 112 (2008) 1040.

\end{enumerate}

\subsection*{Bond-order and species analysis}

\begin{enumerate}[resume, label={[\arabic*]}, leftmargin=2.2em]

  \item J.\ Tersoff,
	``New empirical approach for the structure and energy of covalent systems,''
	\textit{Physical Review B}, 37 (1988) 6991.

  \item LAMMPS \texttt{compute bond/local} and \texttt{compute coord/atom}
	documentation. \url{https://docs.lammps.org/compute.html}

\end{enumerate}

\subsection*{MD and materials defect analysis}

\begin{enumerate}[resume, label={[\arabic*]}, leftmargin=2.2em]

  \item A.\ Stukowski,
	``Visualization and analysis of atomistic simulation data with
	OVITO --- the Open Visualization Tool,''
	\textit{Modelling Simul.\ Mater.\ Sci.\ Eng.}, 18 (2010) 015012.

  \item C.L.\ Kelchner, S.J.\ Plimpton, J.C.\ Hamilton,
	``Dislocation nucleation and defect structure during surface indentation,''
	\textit{Physical Review B}, 58 (1998) 11085.

\end{enumerate}

\subsection*{Diffusion and MSD}

\begin{enumerate}[resume, label={[\arabic*]}, leftmargin=2.2em]

  \item D.\ Frenkel, B.\ Smit,
	\textit{Understanding Molecular Simulation}, 2nd ed.\
	Academic Press (2002). Ch.\ 4: Mean-squared displacement.

\end{enumerate}

\subsection*{Theoretical framework}

\begin{enumerate}[resume, label={[\arabic*]}, leftmargin=2.2em]

  \item Inverse Kaluza-Klein Scale Symmetry Framework.
	VSEPR-SIM Internal Document Series.
	\textit{IKK\_Doc\_I\_Schrodinger\_Alternative.tex},
	\textit{IKK\_Doc\_II\_Initial\_Draft.tex}.
	VSEPR-SIM Theoretical Division, Day 64.

\end{enumerate}

\subsection*{Reserved: CERN / high-energy}

\begin{enumerate}[resume, label={[\arabic*]}, leftmargin=2.2em]

  \item CERN Open Data Portal. \url{https://opendata.cern.ch}
	(Reserved for CERN high-energy verification layer.)

  \item CMS Collaboration, ``Observation of a new boson at a mass of
	125 GeV with the CMS experiment at the LHC,''
	\textit{Physics Letters B}, 716 (2012) 30.
	(Reserved for detector-projection validation, CERN layer.)

\end{enumerate}

% ============================================================================
\section*{Future Work}
\addcontentsline{toc}{section}{Future Work}
\label{sec:futurework}
% ============================================================================

Items are listed in approximate priority order.
Dependency notation: item $A \to B$ means $A$ should precede $B$.

\begin{enumerate}[leftmargin=2em]

  \item \textbf{Native LAMMPS compute/fix integration.}
	Port the post-processing analysis layer into a LAMMPS
	\texttt{compute} or \texttt{fix} for in-situ per-step computation.
	Currently the layer reads \texttt{dump} files offline.
	In-situ avoids I/O bottleneck on large systems.
	\textit{Depends on: stable post-processing API.}

  \item \textbf{Live residual-stream output.}
	Stream per-step residual metrics to a JSON-lines or SSE endpoint
	compatible with the VSEPR-SIM dashboard record format.
	Enables real-time monitoring of packing loss and event detection
	during long production runs.
	\textit{Depends on: item 1 or a separate Python bridge daemon.}

  \item \textbf{VSEPR-SIM xyzFull bridge.}
	Extend the \texttt{.xyzfull} format to carry per-atom
	residual columns, fuzz-radius annotations, and event markers.
	Enables the existing VSEPR-SIM visualization layer to render
	residual overlays directly.
	\textit{Depends on: xyzFull format extension WO; Section 9 schema.}

  \item \textbf{Batch validation campaigns.}
	Automate the nine validation cases in Appendix~F using the
	VSEPR-SIM batch system (\texttt{[batch]}, \texttt{[batch.design]}).
	Produces a reproducible validation matrix with hashes and
	pass/fail ledger.
	\textit{Depends on: Section 10 input files; WO-VSIM-62C batch runner.}

  \item \textbf{Solid-state defect benchmark suite.}
	Extend the vacancy/interstitial case to a systematic
	benchmark: vacancy, interstitial, Frenkel pair, Schottky pair,
	dislocation loop, grain boundary.
	Calibrate $\varepsilon_{A \to S}$ threshold bands against
	known defect formation energies.
	\textit{Depends on: item 4.}

  \item \textbf{Reactive chemistry event benchmark suite.}
	Systematic benchmark against known ReaxFF reaction channels for
	hydrocarbons, water splitting, and metal oxidation.
	Calibrate confidence score weights and persistence thresholds.
	\textit{Depends on: item 4; Appendix E schema implementation.}

  \item \textbf{Detector projection / CERN verification layer.}
	Implement $P_{\mathrm{det}}$ and $\mathcal{R}_{\mathrm{det}}$
	using CERN Open Data collision events.
	Begin with $J/\psi \to \mu^+\mu^-$ (clean dilepton channel)
	before moving to $H \to ZZ \to 4\ell$.
	\textit{Depends on: LAMMPS layer validated; items 5 and 6 complete.}

  \item \textbf{High-energy annihilation-map module.}
	Full implementation of the annihilation map
	$\Ann: \Ident{\mathrm{parent}} \to \Comp{\mathrm{final}}$
	and detector reconstruction
	$\Comp{\mathrm{final}} \to \widehat{\Ident{\mathrm{parent}}}$.
	This is the collision-physics analogue of
	inverse packing.
	\textit{Depends on: item 7.}

  \item \textbf{Dark-sector projection toy model.}
	Implement the dark density residual
	$\DarkDens = P_G(\rho) - P_{\mathrm{EM}}(\rho)$
	as a numerical toy model using parameterised
	kernels $K_G(w)$ and $K_{\mathrm{EM}}(w)$.
	Goal: derive constraints on $\chi_w(w)$ from
	known dark-matter fraction $f_{\mathrm{dark}} \approx 0.27$.
	\textit{Depends on: items 7 and 8; IKK theory document.}

  \item \textbf{Full theory-to-simulation companion paper.}
	Joint document connecting the IKK theoretical framework
	(Documents I and II) to the validated LAMMPS implementation.
	Includes derivation of projection operators from the
	scale ladder, residual-metric physical interpretation,
	and quantitative comparison with existing MD literature.
	\textit{Depends on: all validation campaigns complete.}

\end{enumerate}

% ============================================================================
\vfill
\begin{center}
  \rule{0.4\textwidth}{0.4pt}\\[0.5em]
  \textit{Scale-Packing MD Implementation ---
	VSEPR-SIM-ORGANICS \& METALLICS \{release\}\\
	Back Matter complete.
	Verification is the focus.}
\end{center}

\end{document}
